% Options for packages loaded elsewhere
\PassOptionsToPackage{unicode}{hyperref}
\PassOptionsToPackage{hyphens}{url}
\documentclass[
]{article}
\usepackage{xcolor}
\usepackage{amsmath,amssymb}
\setcounter{secnumdepth}{-\maxdimen} % remove section numbering
\usepackage{iftex}
\ifPDFTeX
  \usepackage[T1]{fontenc}
  \usepackage[utf8]{inputenc}
  \usepackage{textcomp} % provide euro and other symbols
\else % if luatex or xetex
  \usepackage{unicode-math} % this also loads fontspec
  \defaultfontfeatures{Scale=MatchLowercase}
  \defaultfontfeatures[\rmfamily]{Ligatures=TeX,Scale=1}
\fi
\usepackage{lmodern}
\ifPDFTeX\else
  % xetex/luatex font selection
\fi
% Use upquote if available, for straight quotes in verbatim environments
\IfFileExists{upquote.sty}{\usepackage{upquote}}{}
\IfFileExists{microtype.sty}{% use microtype if available
  \usepackage[]{microtype}
  \UseMicrotypeSet[protrusion]{basicmath} % disable protrusion for tt fonts
}{}
\makeatletter
\@ifundefined{KOMAClassName}{% if non-KOMA class
  \IfFileExists{parskip.sty}{%
    \usepackage{parskip}
  }{% else
    \setlength{\parindent}{0pt}
    \setlength{\parskip}{6pt plus 2pt minus 1pt}}
}{% if KOMA class
  \KOMAoptions{parskip=half}}
\makeatother
\usepackage{color}
\usepackage{fancyvrb}
\newcommand{\VerbBar}{|}
\newcommand{\VERB}{\Verb[commandchars=\\\{\}]}
\DefineVerbatimEnvironment{Highlighting}{Verbatim}{commandchars=\\\{\}}
% Add ',fontsize=\small' for more characters per line
\newenvironment{Shaded}{}{}
\newcommand{\AlertTok}[1]{\textcolor[rgb]{1.00,0.00,0.00}{\textbf{#1}}}
\newcommand{\AnnotationTok}[1]{\textcolor[rgb]{0.38,0.63,0.69}{\textbf{\textit{#1}}}}
\newcommand{\AttributeTok}[1]{\textcolor[rgb]{0.49,0.56,0.16}{#1}}
\newcommand{\BaseNTok}[1]{\textcolor[rgb]{0.25,0.63,0.44}{#1}}
\newcommand{\BuiltInTok}[1]{\textcolor[rgb]{0.00,0.50,0.00}{#1}}
\newcommand{\CharTok}[1]{\textcolor[rgb]{0.25,0.44,0.63}{#1}}
\newcommand{\CommentTok}[1]{\textcolor[rgb]{0.38,0.63,0.69}{\textit{#1}}}
\newcommand{\CommentVarTok}[1]{\textcolor[rgb]{0.38,0.63,0.69}{\textbf{\textit{#1}}}}
\newcommand{\ConstantTok}[1]{\textcolor[rgb]{0.53,0.00,0.00}{#1}}
\newcommand{\ControlFlowTok}[1]{\textcolor[rgb]{0.00,0.44,0.13}{\textbf{#1}}}
\newcommand{\DataTypeTok}[1]{\textcolor[rgb]{0.56,0.13,0.00}{#1}}
\newcommand{\DecValTok}[1]{\textcolor[rgb]{0.25,0.63,0.44}{#1}}
\newcommand{\DocumentationTok}[1]{\textcolor[rgb]{0.73,0.13,0.13}{\textit{#1}}}
\newcommand{\ErrorTok}[1]{\textcolor[rgb]{1.00,0.00,0.00}{\textbf{#1}}}
\newcommand{\ExtensionTok}[1]{#1}
\newcommand{\FloatTok}[1]{\textcolor[rgb]{0.25,0.63,0.44}{#1}}
\newcommand{\FunctionTok}[1]{\textcolor[rgb]{0.02,0.16,0.49}{#1}}
\newcommand{\ImportTok}[1]{\textcolor[rgb]{0.00,0.50,0.00}{\textbf{#1}}}
\newcommand{\InformationTok}[1]{\textcolor[rgb]{0.38,0.63,0.69}{\textbf{\textit{#1}}}}
\newcommand{\KeywordTok}[1]{\textcolor[rgb]{0.00,0.44,0.13}{\textbf{#1}}}
\newcommand{\NormalTok}[1]{#1}
\newcommand{\OperatorTok}[1]{\textcolor[rgb]{0.40,0.40,0.40}{#1}}
\newcommand{\OtherTok}[1]{\textcolor[rgb]{0.00,0.44,0.13}{#1}}
\newcommand{\PreprocessorTok}[1]{\textcolor[rgb]{0.74,0.48,0.00}{#1}}
\newcommand{\RegionMarkerTok}[1]{#1}
\newcommand{\SpecialCharTok}[1]{\textcolor[rgb]{0.25,0.44,0.63}{#1}}
\newcommand{\SpecialStringTok}[1]{\textcolor[rgb]{0.73,0.40,0.53}{#1}}
\newcommand{\StringTok}[1]{\textcolor[rgb]{0.25,0.44,0.63}{#1}}
\newcommand{\VariableTok}[1]{\textcolor[rgb]{0.10,0.09,0.49}{#1}}
\newcommand{\VerbatimStringTok}[1]{\textcolor[rgb]{0.25,0.44,0.63}{#1}}
\newcommand{\WarningTok}[1]{\textcolor[rgb]{0.38,0.63,0.69}{\textbf{\textit{#1}}}}
\usepackage{longtable,booktabs,array}
\newcounter{none} % for unnumbered tables
\usepackage{calc} % for calculating minipage widths
% Correct order of tables after \paragraph or \subparagraph
\usepackage{etoolbox}
\makeatletter
\patchcmd\longtable{\par}{\if@noskipsec\mbox{}\fi\par}{}{}
\makeatother
% Allow footnotes in longtable head/foot
\IfFileExists{footnotehyper.sty}{\usepackage{footnotehyper}}{\usepackage{footnote}}
\makesavenoteenv{longtable}
\setlength{\emergencystretch}{3em} % prevent overfull lines
\providecommand{\tightlist}{%
  \setlength{\itemsep}{0pt}\setlength{\parskip}{0pt}}
\usepackage{bookmark}
\IfFileExists{xurl.sty}{\usepackage{xurl}}{} % add URL line breaks if available
\urlstyle{same}
\hypersetup{
  hidelinks,
  pdfcreator={LaTeX via pandoc}}

\author{}
\date{}

\begin{document}

\section{Lux Database Layer Audit
Report}\label{lux-database-layer-audit-report}

\textbf{Date}: 2025-12-30 \textbf{Auditor}: Claude Code (Opus 4.5)
\textbf{Scope}: \texttt{/Users/z/work/lux/node/internal/database/},
\texttt{/Users/z/work/lux/node/x/merkledb/},
\texttt{/Users/z/work/lux/node/x/blockdb/},
\texttt{/Users/z/work/lux/node/x/archivedb/},
\texttt{/Users/z/work/lux/node/cache/}

\begin{center}\rule{0.5\linewidth}{0.5pt}\end{center}

\subsection{Executive Summary}\label{executive-summary}

The Lux database layer demonstrates sound engineering with multiple
layers of crash recovery, checksumming, and concurrent access control.
The architecture separates concerns effectively:

{\def\LTcaptype{none} % do not increment counter
\begin{longtable}[]{@{}lll@{}}
\toprule\noalign{}
Component & Risk Level & Key Finding \\
\midrule\noalign{}
\endhead
\bottomrule\noalign{}
\endlastfoot
BadgerDB Integration & Low & Thin wrapper, delegates to
\texttt{luxfi/database} \\
PebbleDB Support & Low & Alternative backend, same interface \\
MerkleDB & Medium & Complex but well-designed; rebuild on unclean
shutdown \\
BlockDB & Low & Strong crash recovery with checksums \\
ArchiveDB & Low & Simple append-only design \\
Caching & Medium & Size-based LRU; potential memory pressure under
load \\
\end{longtable}
}

\textbf{Overall Assessment}: Production-ready with minor recommendations
for 2025.

\begin{center}\rule{0.5\linewidth}{0.5pt}\end{center}

\subsection{1. BadgerDB Integration}\label{1-badgerdb-integration}

\subsubsection{Implementation}\label{implementation}

\begin{itemize}
\tightlist
\item
  \textbf{Location}:
  \texttt{/Users/z/work/lux/node/internal/database/factory/factory.go}
\item
  \textbf{Design}: Thin wrapper around
  \texttt{github.com/luxfi/database/factory.New}
\end{itemize}

\begin{Shaded}
\begin{Highlighting}[]
\KeywordTok{func}\NormalTok{ New}\OperatorTok{(}\NormalTok{name }\DataTypeTok{string}\OperatorTok{,}\NormalTok{ path }\DataTypeTok{string}\OperatorTok{,}\NormalTok{ readOnly }\DataTypeTok{bool}\OperatorTok{,}\NormalTok{ config }\OperatorTok{[]}\DataTypeTok{byte}\OperatorTok{,}
\NormalTok{    gatherer metrics}\OperatorTok{.}\NormalTok{MultiGatherer}\OperatorTok{,}\NormalTok{ logger log}\OperatorTok{.}\NormalTok{Logger}\OperatorTok{,}
\NormalTok{    metricsPrefix }\DataTypeTok{string}\OperatorTok{,}\NormalTok{ meterDBRegName }\DataTypeTok{string}\OperatorTok{)} \OperatorTok{(}\NormalTok{database}\OperatorTok{.}\NormalTok{Database}\OperatorTok{,} \DataTypeTok{error}\OperatorTok{)} \OperatorTok{\{}
    \ControlFlowTok{return}\NormalTok{ factory}\OperatorTok{.}\NormalTok{New}\OperatorTok{(}\NormalTok{name}\OperatorTok{,}\NormalTok{ path}\OperatorTok{,}\NormalTok{ readOnly}\OperatorTok{,}\NormalTok{ config}\OperatorTok{,}\NormalTok{ gatherer}\OperatorTok{,}\NormalTok{ logger}\OperatorTok{,}\NormalTok{ metricsPrefix}\OperatorTok{,}\NormalTok{ meterDBRegName}\OperatorTok{)}
\OperatorTok{\}}
\end{Highlighting}
\end{Shaded}

\subsubsection{Findings}\label{findings}

\begin{itemize}
\tightlist
\item
  \textbf{Data Corruption Risk}: LOW - BadgerDB handles its own crash
  recovery internally
\item
  \textbf{Batch Capacity}: Constants in \texttt{common.go} suggest
  awareness of write amplification
\end{itemize}

\begin{Shaded}
\begin{Highlighting}[]
\KeywordTok{const} \OperatorTok{(}
\NormalTok{    BatchCapacityCap }\OperatorTok{=} \DecValTok{128} \OperatorTok{*}\NormalTok{ units}\OperatorTok{.}\NormalTok{MiB  }\CommentTok{// Maximum batch size}
\NormalTok{    BatchCapacityMin }\OperatorTok{=} \DecValTok{1} \OperatorTok{*}\NormalTok{ units}\OperatorTok{.}\NormalTok{MiB    }\CommentTok{// Minimum batch size}
\OperatorTok{)}
\end{Highlighting}
\end{Shaded}

\subsubsection{Recommendations}\label{recommendations}

\begin{enumerate}
\def\labelenumi{\arabic{enumi}.}
\tightlist
\item
  Consider exposing BadgerDB\textquotesingle s value log GC tuning for
  long-running nodes
\item
  Add metrics for LSM compaction pressure
\end{enumerate}

\begin{center}\rule{0.5\linewidth}{0.5pt}\end{center}

\subsection{2. PebbleDB Support}\label{2-pebbledb-support}

\subsubsection{Implementation}\label{implementation-1}

\begin{itemize}
\tightlist
\item
  Available as \texttt{"pebbledb"} option in factory
\item
  Same \texttt{database.Database} interface as BadgerDB
\end{itemize}

\subsubsection{Findings}\label{findings-1}

\begin{itemize}
\tightlist
\item
  \textbf{Alternative Backend}: Provides choice for operators
\item
  \textbf{Interface Compliance}: Uses standard
  \texttt{database.Database} interface
\end{itemize}

\subsubsection{Recommendations}\label{recommendations-1}

\begin{enumerate}
\def\labelenumi{\arabic{enumi}.}
\tightlist
\item
  Document performance characteristics vs BadgerDB for operator guidance
\item
  Consider making PebbleDB the default (RocksDB lineage, battle-tested)
\end{enumerate}

\begin{center}\rule{0.5\linewidth}{0.5pt}\end{center}

\subsection{3. MerkleDB}\label{3-merkledb}

\subsubsection{Implementation}\label{implementation-2}

\begin{itemize}
\tightlist
\item
  \textbf{Location}: \texttt{/Users/z/work/lux/node/x/merkledb/}
\item
  \textbf{Design}: Radix trie with Merkle proofs, MVCC views
\end{itemize}

\subsubsection{Architecture}\label{architecture}

\begin{verbatim}
                    ┌─────────────────┐
                    │     merkleDB    │
                    │  (db.go:1400+)  │
                    └────────┬────────┘
                             │
        ┌────────────────────┼────────────────────┐
        │                    │                    │
┌───────▼───────┐   ┌────────▼────────┐   ┌───────▼───────┐
│intermediateDB │   │   valueNodeDB   │   │   onEvictCache│
│ (write buffer)│   │   (LRU cache)   │   │   (FIFO)      │
└───────────────┘   └─────────────────┘   └───────────────┘
\end{verbatim}

\subsubsection{Crash Recovery Mechanism}\label{crash-recovery-mechanism}

\begin{Shaded}
\begin{Highlighting}[]
\CommentTok{// Marker key for clean shutdown detection}
\NormalTok{cleanShutdownKey }\OperatorTok{=} \OperatorTok{[]}\DataTypeTok{byte}\OperatorTok{(}\DataTypeTok{string}\OperatorTok{(}\NormalTok{metadataPrefix}\OperatorTok{)} \OperatorTok{+} \StringTok{"cleanShutdown"}\OperatorTok{)}
\NormalTok{hadCleanShutdown }\OperatorTok{=} \OperatorTok{[]}\DataTypeTok{byte}\OperatorTok{\{}\DecValTok{1}\OperatorTok{\}}
\NormalTok{didNotHaveCleanShutdown }\OperatorTok{=} \OperatorTok{[]}\DataTypeTok{byte}\OperatorTok{\{}\DecValTok{0}\OperatorTok{\}}
\end{Highlighting}
\end{Shaded}

\textbf{Recovery Flow}:

\begin{enumerate}
\def\labelenumi{\arabic{enumi}.}
\tightlist
\item
  On open: Check \texttt{cleanShutdownKey} value
\item
  If \texttt{didNotHaveCleanShutdown}: Call \texttt{rebuild()} to
  reconstruct trie
\item
  Immediately write \texttt{didNotHaveCleanShutdown} marker
\item
  On clean close: Write \texttt{hadCleanShutdown} marker
\end{enumerate}

\subsubsection{Concurrency Model}\label{concurrency-model}

\textbf{Dual Lock Pattern}:

\begin{Shaded}
\begin{Highlighting}[]
\KeywordTok{type}\NormalTok{ merkleDB }\KeywordTok{struct} \OperatorTok{\{}
\NormalTok{    lock       sync}\OperatorTok{.}\NormalTok{RWMutex  }\CommentTok{// Protects general state}
\NormalTok{    commitLock sync}\OperatorTok{.}\NormalTok{RWMutex  }\CommentTok{// Prevents concurrent commits}
    \CommentTok{// ...}
\OperatorTok{\}}
\end{Highlighting}
\end{Shaded}

\textbf{View-Based MVCC}:

\begin{itemize}
\tightlist
\item
  \texttt{NewView()} creates isolated snapshots
\item
  Views share read access, commits serialize via \texttt{commitLock}
\end{itemize}

\subsubsection{Findings}\label{findings-2}

{\def\LTcaptype{none} % do not increment counter
\begin{longtable}[]{@{}lll@{}}
\toprule\noalign{}
Area & Risk & Details \\
\midrule\noalign{}
\endhead
\bottomrule\noalign{}
\endlastfoot
Data Corruption & LOW & Rebuilds from raw key-values on crash \\
Race Conditions & LOW & Proper dual-lock prevents concurrent commits \\
Memory Usage & MEDIUM & Unbounded view accumulation possible \\
Performance & GOOD & Write buffers, batch eviction \\
\end{longtable}
}

\subsubsection{Critical Code Paths}\label{critical-code-paths}

\textbf{Write Buffer Eviction} (\texttt{intermediate\_node\_db.go}):

\begin{Shaded}
\begin{Highlighting}[]
\CommentTok{// Fatal error handling {-} closes database on write failures}
\ControlFlowTok{if}\NormalTok{ err }\OperatorTok{:=}\NormalTok{ i}\OperatorTok{.}\NormalTok{evictToDisk}\OperatorTok{(}\NormalTok{evictionBatchSize}\OperatorTok{);}\NormalTok{ err }\OperatorTok{!=} \OtherTok{nil} \OperatorTok{\{}
    \ControlFlowTok{return} \OtherTok{nil}\OperatorTok{,}\NormalTok{ err  }\CommentTok{// Propagates up, triggers shutdown}
\OperatorTok{\}}
\end{Highlighting}
\end{Shaded}

\textbf{Atomic Closed State} (\texttt{value\_node\_db.go}):

\begin{Shaded}
\begin{Highlighting}[]
\KeywordTok{type}\NormalTok{ valueNodeDB }\KeywordTok{struct} \OperatorTok{\{}
\NormalTok{    closed atomic}\OperatorTok{.}\NormalTok{Bool  }\CommentTok{// Safe concurrent close detection}
    \CommentTok{// ...}
\OperatorTok{\}}
\end{Highlighting}
\end{Shaded}

\subsubsection{Recommendations}\label{recommendations-2}

\begin{enumerate}
\def\labelenumi{\arabic{enumi}.}
\tightlist
\item
  Add view lifecycle limits to prevent memory exhaustion
\item
  Consider background compaction for intermediate node DB
\item
  Add metrics for rebuild time on recovery
\end{enumerate}

\begin{center}\rule{0.5\linewidth}{0.5pt}\end{center}

\subsection{4. BlockDB}\label{4-blockdb}

\subsubsection{Implementation}\label{implementation-3}

\begin{itemize}
\tightlist
\item
  \textbf{Location}:
  \texttt{/Users/z/work/lux/node/x/blockdb/database.go}
\item
  \textbf{Design}: Block storage with index file and data files
\end{itemize}

\subsubsection{Crash Recovery
Mechanism}\label{crash-recovery-mechanism-1}

\textbf{Index File Structure}:

\begin{verbatim}
┌──────────────────┬──────────────────┬─────────────────────┐
│ Version Header   │ Block Index      │ Checksum            │
│ (magic + version)│ (height->offset) │ (xxhash per block)  │
└──────────────────┴──────────────────┴─────────────────────┘
\end{verbatim}

\textbf{Recovery Algorithm}:

\begin{enumerate}
\def\labelenumi{\arabic{enumi}.}
\tightlist
\item
  Validate index file header
\item
  If corrupted: Scan data files sequentially
\item
  Verify each block\textquotesingle s xxhash checksum
\item
  Rebuild index from valid blocks
\item
  Truncate at first corrupted block
\end{enumerate}

\subsubsection{Concurrency Model}\label{concurrency-model-1}

\begin{Shaded}
\begin{Highlighting}[]
\KeywordTok{type}\NormalTok{ Database }\KeywordTok{struct} \OperatorTok{\{}
\NormalTok{    closeMu             sync}\OperatorTok{.}\NormalTok{RWMutex           }\CommentTok{// Close coordination}
\NormalTok{    fileOpenMu          sync}\OperatorTok{.}\NormalTok{Mutex             }\CommentTok{// File handle management}
\NormalTok{    blockHeights        atomic}\OperatorTok{.}\NormalTok{Pointer}\OperatorTok{[...]}    \CommentTok{// Lock{-}free height queries}
\NormalTok{    nextDataWriteOffset atomic}\OperatorTok{.}\NormalTok{Uint64          }\CommentTok{// Lock{-}free offset tracking}
\NormalTok{    headerWriteOccupied atomic}\OperatorTok{.}\NormalTok{Bool            }\CommentTok{// Single header writer}
\OperatorTok{\}}
\end{Highlighting}
\end{Shaded}

\textbf{Atomic Operations}:

\begin{itemize}
\tightlist
\item
  \texttt{blockHeights} uses atomic pointer swap for lock-free reads
\item
  \texttt{nextDataWriteOffset} enables concurrent write offset
  calculation
\item
  \texttt{headerWriteOccupied} prevents concurrent header updates
\end{itemize}

\subsubsection{Compression}\label{compression}

\begin{itemize}
\tightlist
\item
  Uses zstd compression for block data
\item
  Configurable compression level
\end{itemize}

\subsubsection{Findings}\label{findings-3}

{\def\LTcaptype{none} % do not increment counter
\begin{longtable}[]{@{}lll@{}}
\toprule\noalign{}
Area & Risk & Details \\
\midrule\noalign{}
\endhead
\bottomrule\noalign{}
\endlastfoot
Data Corruption & VERY LOW & xxhash checksums, recovery scan \\
Race Conditions & LOW & Well-designed atomic/mutex hybrid \\
Memory Usage & LOW & Streaming I/O, minimal buffering \\
Performance & GOOD & Lock-free reads, batch writes \\
\end{longtable}
}

\subsubsection{Test Coverage}\label{test-coverage}

\begin{itemize}
\tightlist
\item
  \texttt{recovery\_test.go} covers various corruption scenarios
\item
  Tests include partial writes, truncated files, corrupted checksums
\end{itemize}

\subsubsection{Recommendations}\label{recommendations-3}

\begin{enumerate}
\def\labelenumi{\arabic{enumi}.}
\tightlist
\item
  Add periodic background checksum verification
\item
  Consider index file replication for faster recovery
\end{enumerate}

\begin{center}\rule{0.5\linewidth}{0.5pt}\end{center}

\subsection{5. ArchiveDB}\label{5-archivedb}

\subsubsection{Implementation}\label{implementation-4}

\begin{itemize}
\tightlist
\item
  \textbf{Location}: \texttt{/Users/z/work/lux/node/x/archivedb/db.go}
\item
  \textbf{Design}: Append-only height-indexed state changes
\end{itemize}

\subsubsection{Findings}\label{findings-4}

\begin{itemize}
\tightlist
\item
  \textbf{Simplest Component}: Thin wrapper for height-based queries
\item
  \textbf{No Complex Recovery}: Relies on underlying
  database\textquotesingle s durability
\item
  \textbf{Risk Level}: LOW
\end{itemize}

\begin{center}\rule{0.5\linewidth}{0.5pt}\end{center}

\subsection{6. Caching Layer}\label{6-caching-layer}

\subsubsection{Implementation}\label{implementation-5}

\begin{itemize}
\tightlist
\item
  \textbf{Location}: \texttt{/Users/z/work/lux/node/cache/}
\end{itemize}

\subsubsection{LRU Cache Types}\label{lru-cache-types}

{\def\LTcaptype{none} % do not increment counter
\begin{longtable}[]{@{}lll@{}}
\toprule\noalign{}
Type & Location & Use Case \\
\midrule\noalign{}
\endhead
\bottomrule\noalign{}
\endlastfoot
\texttt{SizedLRU} & \texttt{lru/sized\_cache.go} & Size-based
eviction \\
\texttt{LRU} & \texttt{lru\_cache.go} & Count-based eviction \\
\texttt{onEvictCache} & \texttt{merkledb/cache.go} & Callback on
eviction \\
\end{longtable}
}

\subsubsection{SizedLRU Implementation}\label{sizedlru-implementation}

\begin{Shaded}
\begin{Highlighting}[]
\KeywordTok{type}\NormalTok{ SizedLRU}\OperatorTok{[}\NormalTok{K }\DataTypeTok{comparable}\OperatorTok{,}\NormalTok{ V }\DataTypeTok{any}\OperatorTok{]} \KeywordTok{struct} \OperatorTok{\{}
\NormalTok{    lock        sync}\OperatorTok{.}\NormalTok{Mutex}
\NormalTok{    elements    }\OperatorTok{*}\NormalTok{linked}\OperatorTok{.}\NormalTok{Hashmap}\OperatorTok{[}\NormalTok{K}\OperatorTok{,} \OperatorTok{*}\NormalTok{sizedElement}\OperatorTok{[}\NormalTok{V}\OperatorTok{]]}
\NormalTok{    maxSize     }\DataTypeTok{int}
\NormalTok{    currentSize }\DataTypeTok{int}
\OperatorTok{\}}

\KeywordTok{type}\NormalTok{ sizedElement}\OperatorTok{[}\NormalTok{V }\DataTypeTok{any}\OperatorTok{]} \KeywordTok{struct} \OperatorTok{\{}
\NormalTok{    Value V}
\NormalTok{    Size  }\DataTypeTok{int}  \CommentTok{// Stored size prevents recalculation inconsistencies}
\OperatorTok{\}}
\end{Highlighting}
\end{Shaded}

\textbf{Key Design Decision}: Size stored with element prevents
inconsistencies if value is modified after insertion.

\subsubsection{onEvictCache (MerkleDB)}\label{onevictcache-merkledb}

\begin{Shaded}
\begin{Highlighting}[]
\KeywordTok{type}\NormalTok{ onEvictCache}\OperatorTok{[}\NormalTok{K }\DataTypeTok{comparable}\OperatorTok{,}\NormalTok{ V }\DataTypeTok{any}\OperatorTok{]} \KeywordTok{struct} \OperatorTok{\{}
\NormalTok{    lock       sync}\OperatorTok{.}\NormalTok{RWMutex}
\NormalTok{    maxSize    }\DataTypeTok{int}
\NormalTok{    fifo       }\OperatorTok{*}\NormalTok{linked}\OperatorTok{.}\NormalTok{Hashmap}\OperatorTok{[}\NormalTok{K}\OperatorTok{,}\NormalTok{ V}\OperatorTok{]}  \CommentTok{// FIFO for eviction order}
\NormalTok{    onEviction }\KeywordTok{func}\OperatorTok{(}\NormalTok{K}\OperatorTok{,}\NormalTok{ V}\OperatorTok{)} \DataTypeTok{error}       \CommentTok{// Callback on evict}
\OperatorTok{\}}
\end{Highlighting}
\end{Shaded}

\textbf{FIFO vs LRU}: Uses FIFO for eviction order, not LRU. This is
intentional for write buffer semantics.

\subsubsection{Findings}\label{findings-5}

{\def\LTcaptype{none} % do not increment counter
\begin{longtable}[]{@{}lll@{}}
\toprule\noalign{}
Area & Risk & Details \\
\midrule\noalign{}
\endhead
\bottomrule\noalign{}
\endlastfoot
Memory Leaks & LOW & Proper size tracking, eviction callbacks \\
Race Conditions & LOW & Mutex-protected operations \\
Performance & MEDIUM & Lock contention possible under high load \\
\end{longtable}
}

\subsubsection{Recommendations}\label{recommendations-4}

\begin{enumerate}
\def\labelenumi{\arabic{enumi}.}
\tightlist
\item
  Consider sharded caches for reduced lock contention
\item
  Add cache hit/miss metrics for tuning
\item
  Evaluate ARC or 2Q algorithms for better hit rates
\end{enumerate}

\begin{center}\rule{0.5\linewidth}{0.5pt}\end{center}

\subsection{Data Integrity Analysis}\label{data-integrity-analysis}

\subsubsection{Defense in Depth}\label{defense-in-depth}

\begin{verbatim}
Layer 1: Application Checksums (xxhash in BlockDB)
    │
    ▼
Layer 2: Clean Shutdown Markers (MerkleDB)
    │
    ▼
Layer 3: Database Engine Recovery (BadgerDB/PebbleDB)
    │
    ▼
Layer 4: Filesystem Journaling (ext4/APFS)
\end{verbatim}

\subsubsection{Corruption Scenarios
Handled}\label{corruption-scenarios-handled}

{\def\LTcaptype{none} % do not increment counter
\begin{longtable}[]{@{}lll@{}}
\toprule\noalign{}
Scenario & MerkleDB & BlockDB \\
\midrule\noalign{}
\endhead
\bottomrule\noalign{}
\endlastfoot
Process crash mid-write & Rebuild from raw KV & Checksum validation \\
Partial block write & N/A & Truncate at corruption \\
Index corruption & Rebuild trie & Scan data files \\
Power failure & Rebuild from raw KV & Rebuild index \\
\end{longtable}
}

\subsubsection{Gap Analysis}\label{gap-analysis}

{\def\LTcaptype{none} % do not increment counter
\begin{longtable}[]{@{}lll@{}}
\toprule\noalign{}
Scenario & Current Handling & Risk \\
\midrule\noalign{}
\endhead
\bottomrule\noalign{}
\endlastfoot
Bit rot in cold data & Not detected & LOW (rare) \\
Concurrent file access & Mutex protected & LOW \\
Full disk & Propagates error & MEDIUM \\
\end{longtable}
}

\begin{center}\rule{0.5\linewidth}{0.5pt}\end{center}

\subsection{Performance Analysis}\label{performance-analysis}

\subsubsection{Write Path}\label{write-path}

\begin{verbatim}
Application Write
    │
    ▼
┌───────────────┐
│ Write Buffer  │ ◄── Deferred disk writes
│ (in-memory)   │
└───────┬───────┘
        │ Batch eviction
        ▼
┌───────────────┐
│ BadgerDB/     │ ◄── LSM tree writes
│ PebbleDB      │
└───────────────┘
\end{verbatim}

\textbf{Optimizations}:

\begin{itemize}
\tightlist
\item
  Write buffers reduce disk I/O
\item
  Batch eviction amortizes write costs
\item
  Atomic operations reduce lock overhead
\end{itemize}

\subsubsection{Read Path}\label{read-path}

\begin{verbatim}
Application Read
    │
    ▼
┌───────────────┐
│ LRU Cache     │ ◄── O(1) lookup
└───────┬───────┘
        │ Miss
        ▼
┌───────────────┐
│ Write Buffer  │ ◄── Check pending writes
└───────┬───────┘
        │ Miss
        ▼
┌───────────────┐
│ Database      │ ◄── Disk read
└───────────────┘
\end{verbatim}

\subsubsection{Bottlenecks}\label{bottlenecks}

{\def\LTcaptype{none} % do not increment counter
\begin{longtable}[]{@{}lll@{}}
\toprule\noalign{}
Component & Bottleneck & Mitigation \\
\midrule\noalign{}
\endhead
\bottomrule\noalign{}
\endlastfoot
MerkleDB & Commit serialization & Expected (consistency) \\
BlockDB & Single header writer & Batched updates \\
Cache & Lock contention & Sharding (recommended) \\
\end{longtable}
}

\begin{center}\rule{0.5\linewidth}{0.5pt}\end{center}

\subsection{Security Concerns}\label{security-concerns}

\subsubsection{Low Risk}\label{low-risk}

\begin{enumerate}
\def\labelenumi{\arabic{enumi}.}
\tightlist
\item
  \textbf{No External Input}: Database layer receives validated data
  from VM layer
\item
  \textbf{No Network Exposure}: Internal component, not directly
  accessible
\end{enumerate}

\subsubsection{Moderate Risk}\label{moderate-risk}

\begin{enumerate}
\def\labelenumi{\arabic{enumi}.}
\tightlist
\item
  \textbf{Denial of Service}: Unbounded view creation could exhaust
  memory
\item
  \textbf{Resource Exhaustion}: Large batch operations could pressure
  memory
\end{enumerate}

\subsubsection{Recommendations}\label{recommendations-5}

\begin{enumerate}
\def\labelenumi{\arabic{enumi}.}
\tightlist
\item
  Add view count limits per database instance
\item
  Implement memory pressure backoff for batch operations
\end{enumerate}

\begin{center}\rule{0.5\linewidth}{0.5pt}\end{center}

\subsection{2025 Recommendations}\label{2025-recommendations}

\subsubsection{Priority 1:
Observability}\label{priority-1-observability}

\begin{itemize}
\tightlist
\item[$\square$]
  Add Prometheus metrics for cache hit rates
\item[$\square$]
  Add metrics for MerkleDB rebuild time
\item[$\square$]
  Add BlockDB checksum verification stats
\end{itemize}

\subsubsection{Priority 2: Performance}\label{priority-2-performance}

\begin{itemize}
\tightlist
\item[$\square$]
  Evaluate sharded caches for high-throughput scenarios
\item[$\square$]
  Consider ARC/2Q cache replacement algorithms
\item[$\square$]
  Profile lock contention under production load
\end{itemize}

\subsubsection{Priority 3: Resilience}\label{priority-3-resilience}

\begin{itemize}
\tightlist
\item[$\square$]
  Implement periodic background checksum verification for BlockDB
\item[$\square$]
  Add memory pressure detection and backoff
\item[$\square$]
  Consider index file replication for faster BlockDB recovery
\end{itemize}

\subsubsection{Priority 4:
Documentation}\label{priority-4-documentation}

\begin{itemize}
\tightlist
\item[$\square$]
  Document BadgerDB vs PebbleDB performance characteristics
\item[$\square$]
  Add runbook for crash recovery scenarios
\item[$\square$]
  Document tuning parameters for different hardware profiles
\end{itemize}

\begin{center}\rule{0.5\linewidth}{0.5pt}\end{center}

\subsection{Appendix A: File Inventory}\label{appendix-a-file-inventory}

{\def\LTcaptype{none} % do not increment counter
\begin{longtable}[]{@{}lll@{}}
\toprule\noalign{}
File & Lines & Purpose \\
\midrule\noalign{}
\endhead
\bottomrule\noalign{}
\endlastfoot
\texttt{internal/database/factory/factory.go} & \textasciitilde30 &
Database creation \\
\texttt{internal/database/common.go} & \textasciitilde20 & Batch
constants \\
\texttt{x/merkledb/db.go} & \textasciitilde1400 & Core MerkleDB \\
\texttt{x/merkledb/cache.go} & \textasciitilde150 & onEvictCache \\
\texttt{x/merkledb/intermediate\_node\_db.go} & \textasciitilde250 &
Write buffer \\
\texttt{x/merkledb/value\_node\_db.go} & \textasciitilde200 & Value
storage \\
\texttt{x/blockdb/database.go} & \textasciitilde1200 & Block storage \\
\texttt{x/archivedb/db.go} & \textasciitilde200 & Archive storage \\
\texttt{cache/lru/sized\_cache.go} & \textasciitilde150 & Size-based
LRU \\
\texttt{cache/lru\_cache.go} & \textasciitilde100 & Basic LRU \\
\end{longtable}
}

\subsection{Appendix B: Key Constants}\label{appendix-b-key-constants}

\begin{Shaded}
\begin{Highlighting}[]
\CommentTok{// Batch capacity management}
\NormalTok{BatchCapacityCap }\OperatorTok{=} \DecValTok{128} \OperatorTok{*}\NormalTok{ units}\OperatorTok{.}\NormalTok{MiB}
\NormalTok{BatchCapacityMin }\OperatorTok{=} \DecValTok{1} \OperatorTok{*}\NormalTok{ units}\OperatorTok{.}\NormalTok{MiB}

\CommentTok{// MerkleDB markers}
\NormalTok{cleanShutdownKey }\OperatorTok{=} \OperatorTok{[]}\DataTypeTok{byte}\OperatorTok{(}\NormalTok{metadataPrefix }\OperatorTok{+} \StringTok{"cleanShutdown"}\OperatorTok{)}
\NormalTok{hadCleanShutdown }\OperatorTok{=} \OperatorTok{[]}\DataTypeTok{byte}\OperatorTok{\{}\DecValTok{1}\OperatorTok{\}}
\NormalTok{didNotHaveCleanShutdown }\OperatorTok{=} \OperatorTok{[]}\DataTypeTok{byte}\OperatorTok{\{}\DecValTok{0}\OperatorTok{\}}

\CommentTok{// BlockDB}
\NormalTok{indexFileVersion }\OperatorTok{=} \DecValTok{1}
\NormalTok{dataFileVersion }\OperatorTok{=} \DecValTok{1}
\end{Highlighting}
\end{Shaded}

\begin{center}\rule{0.5\linewidth}{0.5pt}\end{center}

\textbf{End of Audit Report}

\end{document}
